\section{材料与方法}

\subsection{数据集与评估协议}
本文采用 DUT Anti-UAV 可见光检测子集评估远距离单类别无人机检测性能。官方划分包含 5200 张训练图像、2600 张验证图像和 2200 张测试图像。训练与评估均保持原始长宽比，并采用 letterbox 将输入填充至 $640\times640$ 像素。模型仅依据验证集 \mapall{} 选择 checkpoint，随后在测试集上评估所选模型。

图像级 SHA-256 校验发现训练集与测试集之间存在一张重复图像。为保证数据可追溯性，源文件保持不变，并通过固定评估清单排除该测试副本。最终测试集包含 2199 张图像和 2244 个无人机实例。尺度分析采用与检测器一致的 letterbox 坐标：目标最长边小于 16 像素时记为 tiny，16--32 像素时记为 small，不小于 32 像素时记为 large。这些区间对应 DUT Anti-UAV 中无人机的实际表观尺度。

\subsection{总体网络结构}
SRPA-YOLO 采用单阶段密集检测流程，由多尺度骨干网络、轻量双向颈部和步长分别为 4、8、16 的三个预测层组成。图~\ref{fig:network} 展示完整特征流及训练检测头到部署检测头的变换过程。

\begin{figure*}[!t]
\centering
\includegraphics[width=\textwidth]{Fig1.pdf}
\caption{SRPA-YOLO 总体结构。多尺度骨干与双向颈部输出 P2--P4 特征；训练阶段分别优化共享表征和分类私有表征，精确融合后得到单路径静态部署结构。\label{fig:network}}
\end{figure*}

\subsubsection{多尺度空间特征提取}
设输入图像为 $X_0\in\mathbb{R}^{3\times640\times640}$，骨干第 $l$ 阶段输出可表示为
\begin{equation}
X_l=\mathcal{B}_l(X_{l-1})=\mathcal{A}_l\!\left(\operatorname{BN}(K_l*X_{l-1}+b_l)\right),
\label{eq:backbone}
\end{equation}
其中，$K_l$ 表示空间卷积，$\operatorname{BN}$ 表示批归一化，$\mathcal{A}_l$ 包含 SiLU 激活和局部 C2f 聚合。P2、P3、P4 与 P5 特征的步长为 $s_l\in\{4,8,16,32\}$。对于变换后最长边为 $d$ 像素的目标，其在第 $l$ 层上的近似覆盖范围为
\begin{equation}
n_l=\frac{d}{s_l}.
\label{eq:targetsupport}
\end{equation}
因此，$d=8$ 像素的无人机在 P2 上约覆盖两个特征单元，而在步长 8 的特征层上仅覆盖一个单元。保留 P2 能够提高极小目标的空间采样密度，降低目标在预测前被下采样消除的概率。

\subsubsection{轻量双向特征融合}
颈部同时执行自顶向下的语义传播和自底向上的空间聚合。相邻特征层的自顶向下特征 $T_l$ 与自底向上特征 $N_l$ 表示为
\begin{equation}
\begin{aligned}
T_l&=\psi_l\!\left(\left[X_l,\operatorname{Up}(T_{l+1})\right]\right),\\
N_l&=\eta_l\!\left(\left[T_l,\operatorname{Down}(N_{l-1})\right]\right),
\end{aligned}
\label{eq:neck}
\end{equation}
其中，$[\cdot,\cdot]$ 表示通道拼接，$\operatorname{Up}$ 为最近邻插值，$\operatorname{Down}$ 为深度步长卷积，$\psi_l$ 和 $\eta_l$ 为紧凑的 C2f/Ghost 特征混合器。自顶向下路径将深层语义传递到高分辨率特征，自底向上路径向深层补充边缘和位移信息。对于外观响应较弱、但仍具有局部空间对比的无人机，双向交换有助于同时保持语义判别性与定位敏感性。

\subsubsection{上下文感知特征校准}
高分辨率阶段执行通道与空间校准，深层阶段聚合更宽的上下文。对特征 $F$，校准响应可表示为
\begin{equation}
\begin{aligned}
\widetilde{F}&=F\odot M_c(F)\odot M_s(F),\\
M_c(F)&=\sigma\!\left(g_c(\operatorname{AvgPool}F,\operatorname{MaxPool}F)\right),\\
M_s(F)&=\sigma\!\left(g_s(\operatorname{Avg}_cF,\operatorname{Max}_cF)\right),
\end{aligned}
\label{eq:refinement}
\end{equation}
其中，$\odot$ 表示逐元素乘法，$\sigma$ 为 Sigmoid 函数，$g_c$ 和 $g_s$ 分别为轻量通道与空间变换。P2/P3 上的空间项用于抑制云层边缘、植被和建筑纹理等类无人机响应；深层上下文则扩大有效感受野，使孤立飞行目标与重复局部纹理更易区分。这些模块构成本文保留的基础特征提取器，本文的主要改进集中于预测头内部的任务容量分配。

\subsubsection{P2--P4 预测接口}
双向颈部在 $l\in\{2,3,4\}$ 输出 $N_l$，尺度投影将其映射到检测头输入空间：
\begin{equation}
\begin{aligned}
F_l&=P_l*N_l,\quad F_l\in\mathbb{R}^{C_l\times H_l\times W_l},\\
(H_l,W_l)&\in\{(160,160),(80,80),(40,40)\}.
\end{aligned}
\label{eq:headinput}
\end{equation}
三个尺度为 tiny 和 small 无人机提供密集覆盖，同时避免增加运行时 P5 预测分支。各尺度独立采用相同的 SRPA 分配方式，无需动态路由或尺度相关控制流。

\subsection{共享表征与私有秩分配}
传统共享检测头使分类与定位使用同一特征字典。SRPA 保留两任务共用的共享表征，并仅向分类分配低秩私有残差，如图~\ref{fig:allocation} 所示。

\begin{figure*}[!t]
\centering
\includegraphics[width=0.92\textwidth]{Fig2.pdf}
\caption{传统共享检测头与 SRPA 检测头的任务特征分配。私有表征只参与分类，其到边界框输出的雅可比恒为零。\label{fig:allocation}}
\end{figure*}

对输入特征 $F_l$，共享 RepConv 与分类私有 RepConv 分别产生
\begin{equation}
\begin{aligned}
H_l^{s}&=\phi_s^{(l)}(F_l)\in\mathbb{R}^{64\times H_l\times W_l},\\
H_l^{p}&=\phi_p^{(l)}(F_l)\in\mathbb{R}^{16\times H_l\times W_l}.
\end{aligned}
\label{eq:features}
\end{equation}
设 $z_l^{box}$ 包含 $4K$ 个 DFL logits，且每个坐标使用 $K=16$ 个离散区间；$z_l^{cls}$ 表示无人机类别 logit。SRPA 的输出为
\begin{equation}
z_l^{box}=W_l^{b}H_l^{s}+b_l^{b},
\label{eq:box}
\end{equation}
\begin{equation}
z_l^{cls}=W_l^{s}H_l^{s}+b_l^{s}+W_l^{p}H_l^{p}.
\label{eq:cls}
\end{equation}
其中，64 和 16 通道是实现层面的表征维数，以下分别记为 $r_s$ 和 $r_p$。它们规定共享与私有通道预算，并不假设学习到的卷积矩阵必然达到满代数秩。由式~\eqref{eq:box} 和式~\eqref{eq:cls} 可得
\begin{equation}
\frac{\partial z_l^{box}}{\partial H_l^{p}}=0,
\label{eq:jacobian}
\end{equation}
该关系表明边界框输出在计算图上直接独立于分类私有表征，上游任务交互则由阶段 I 学习到的共享表征承载。将空间维展平后，$H_l^s\in\mathbb{R}^{r_s\times N_l}$、$H_l^p\in\mathbb{R}^{r_p\times N_l}$，其中 $N_l=H_lW_l$。分类 logit 的行空间满足
\begin{equation}
\begin{aligned}
z_l^{cls}-b_l^{cls}\mathbf{1}^{\top}
&\in \operatorname{row}(H_l^s)+\operatorname{row}(H_l^p),\\
\dim\!\left(\operatorname{row}(H_l^s)+\operatorname{row}(H_l^p)\right)
&\le r_s+r_p.
\end{aligned}
\label{eq:rowspace}
\end{equation}
式~\eqref{eq:rowspace} 给出组合行空间维数的上界，实际维数由各分支的学习秩及其子空间重叠共同决定。该式说明分类最多可访问 $r_s+r_p=80$ 个表征维度，而定位仅由 $r_s=64$ 个共享维度计算。私有分类投影采用 $W_l^p=0$ 初始化，并保持插入时的共享检测器不变，因此满足函数保持条件
\begin{equation}
f_{\mathrm{SRPA}}(x;W^p=0)=f_s(x),\qquad \forall x.
\label{eq:identity}
\end{equation}
扩展后的检测器以与已训练共享检测器完全一致的预测函数开始细化。私有细化期间固定骨干、颈部、共享 stem、边界框投影和共享分类投影，仅优化 $\phi_p$ 与 $W^p$。该约束使阶段 II 的既有定位映射保持不变，同时为低对比度无人机提供附加分类容量。

\FloatBarrier

\subsection{精确结构重参数化}
训练阶段的每个 RepConv stem 包含线性卷积--归一化分支。卷积核 $K$ 与批归一化参数 $(\gamma,\beta,\mu,\sigma^2)$ 可转换为
\begin{equation}
\widehat{K}=\frac{\gamma}{\sqrt{\sigma^2+\epsilon}}K,\qquad
\widehat{b}=\beta+\frac{\gamma}{\sqrt{\sigma^2+\epsilon}}(b-\mu).
\label{eq:bnfusion}
\end{equation}
将各分支卷积核补齐为 $3\times3$ 后求和。设共享和私有 stem 的等效参数分别为 $(K_s,b_s)$ 与 $(K_p,b_p)$，其中 $K_s\in\mathbb{R}^{64\times C\times3\times3}$、$K_p\in\mathbb{R}^{16\times C\times3\times3}$，则通道维精确融合为
\begin{equation}
K_d=\begin{bmatrix}K_s\\K_p\end{bmatrix},\qquad
b_d=\begin{bmatrix}b_s\\b_p\end{bmatrix}.
\label{eq:stemfuse}
\end{equation}
当 DFL 分桶数 $K=16$、类别数 $n_c=1$ 时，对应的 $1\times1$ 投影满足 $W_b\in\mathbb{R}^{4K\times64\times1\times1}$、$W_s\in\mathbb{R}^{n_c\times64\times1\times1}$ 和 $W_p\in\mathbb{R}^{n_c\times16\times1\times1}$。部署投影转换为
\begin{equation}
W_{box}^{d}=\begin{bmatrix}W_b&0\end{bmatrix},\qquad
W_{cls}^{d}=\begin{bmatrix}W_s&W_p\end{bmatrix}.
\label{eq:projfuse}
\end{equation}
对于 $H=[H_s;H_p]$，在 $W_{box}^{d}$ 的私有通道列补零可得 $W_{box}^{d}H=W_bH_s$，拼接分类投影则有 $W_{cls}^{d}H=W_sH_s+W_pH_p$。因此，每个尺度的训练期任务分解可转换为单个 80 通道 $3\times3$ 卷积和一个合并的 $1\times1$ 投影，除浮点舍入误差外计算输出不变。图~\ref{fig:reparam} 展示该过程。

\begin{figure}[!t]
\centering
\includegraphics[width=\columnwidth]{Fig3.pdf}
\caption{SRPA 检测头的结构重参数化。并行训练表征提供任务特定容量，精确卷积核与投影融合产生单路径部署图。\label{fig:reparam}}
\end{figure}

图~\ref{fig:srpahead} 进一步展示单尺度检测头及其任务特定投影。

\begin{figure}[!t]
\centering
\includegraphics[width=\columnwidth]{Fig4.pdf}
\caption{SRPA 检测头内部结构。共享特征同时支持分类与定位，16 通道私有表征在部署融合前提供加性分类残差。\label{fig:srpahead}}
\end{figure}

\FloatBarrier

\subsection{目标函数与两阶段优化}
监督目标为
\begin{equation}
\begin{aligned}
\mathcal{L}_{sup}={}&\lambda_{box}\mathcal{L}_{CIoU}
+\lambda_{cls}\mathcal{L}_{BCE}
+\lambda_{dfl}\mathcal{L}_{DFL},\\
(\lambda_{box},\!\lambda_{cls},\!\lambda_{dfl})={}&(7.5,0.5,1.5),
\end{aligned}
\label{eq:sup}
\end{equation}
其中，CIoU 联合重叠率、中心距离与宽高比惩罚\cite{rezatofighi2019giou,zheng2020diou}，BCE 监督类别置信度，DFL 将每个边界框坐标表示为离散分布。

共享检测器采用冻结的 RLRD-YOLO 教师进行正则化。温度 $T=2$ 时，分类响应蒸馏为
\begin{equation}
\begin{aligned}
p_t&=\sigma(z_t^{cls}/T),\qquad p_s=\sigma(z_s^{cls}/T),\\
\mathcal{L}_{clsKD}&=T^2\operatorname{KL}(p_t\,\Vert\,p_s).
\end{aligned}
\label{eq:clskd}
\end{equation}
对正样本位置集合 $\Omega_+$，密集定位蒸馏和置信度聚焦定位蒸馏分别为
\begin{equation}
\mathcal{L}_{dflKD}=\frac{1}{N_+}\sum_{j\in\Omega_+}\operatorname{KL}\!\left(P_t^{(j)}\Vert P_s^{(j)}\right),
\label{eq:dflkd}
\end{equation}
\begin{equation}
\mathcal{L}_{focus}=\frac{\sum_{j\in\Omega_+}q_j^2\operatorname{KL}(P_t^{(j)}\Vert P_s^{(j)})}{\sum_{j\in\Omega_+}q_j^2+\epsilon},
\label{eq:focus}
\end{equation}
其中，$q_j$ 为教师置信度。阶段 I 的目标函数为
\begin{equation}
\mathcal{L}=\mathcal{L}_{sup}+0.25\mathcal{L}_{clsKD}+0.50\mathcal{L}_{dflKD}+0.125\mathcal{L}_{focus}.
\label{eq:total}
\end{equation}
阶段 II 插入零输出分类私有表征，并从式~\eqref{eq:identity} 的函数保持状态仅优化 $\phi_p$ 与 $W^p$。蒸馏过程中教师始终冻结。该训练流程先向紧凑共享检测器传递定位知识，再在阶段 II 边界框函数固定的条件下细化分类置信度。

\subsection{实验环境与训练设置}
详细训练参数见表~\ref{tab:environment}。

\begin{table}[!t]
\caption{实验环境与训练设置。\label{tab:environment}}
\centering\srpatablesetup
\begin{tabular}{@{}p{0.25\columnwidth}p{0.67\columnwidth}@{}}
\toprule
项目 & 设置\\
\midrule
框架 & Ultralytics 8.3.239 / PyTorch 2.9.0\\
CUDA & 12.8\\
GPU & NVIDIA GeForce RTX 5060，8 GB\\
输入尺寸 & $640\times640$\\
批量大小 & 32\\
训练轮数 & 200\\
优化器 & SGD\\
初始学习率 & 0.01\\
动量 & 0.937\\
权重衰减 & $5\times10^{-4}$\\
Workers & 2\\
数据增强 & HSV 抖动、水平翻转、平移、缩放、Mosaic\\
\bottomrule
\end{tabular}
\end{table}

\subsection{评价指标与部署测试}
Precision 与 Recall 定义为
\begin{equation}
P=\frac{TP}{TP+FP},\qquad R=\frac{TP}{TP+FN}.
\label{eq:pr}
\end{equation}
AP@0.50 和 AP@0.75 分别表示 IoU 阈值为 0.50 和 0.75 时 Precision--Recall 曲线下的面积。主要指标为十个阈值上的平均 AP：
\begin{equation}
\mathrm{mAP}_{0.50:0.95}=\frac{1}{10}\sum_{k=0}^{9}\mathrm{AP}_{0.50+0.05k}.
\label{eq:map}
\end{equation}
评估采用 batch size 32、置信度下限 0.001、NMS IoU 0.7，每幅图像最多保留 300 个检测结果。匹配阶段 II 实验采用随机种子 40、41 和 42，各实验使用相同的阶段 I 共享基底，并独立初始化分类私有参数，以评估私有细化过程在不同初始化条件下的性能稳定性。

共享/私有宽度分析中，设配置 $c$ 的验证集 \mapall{} 与 AP@0.75 分别为 $m_c$ 和 $a_c$。首先保留与最高验证 \mapall{} 差值不超过 0.003 的配置：
\begin{equation}
\mathcal{C}_{\epsilon}=\left\{c:m_c\ge \max_j m_j-0.003\right\},
\label{eq:selection}
\end{equation}
再选择 AP@0.75 最高的配置，若完全相同则由 AP@0.50 决定。该规则在保持主要指标的同时优先选择定位更准确的模型。

延迟测试使用 batch 1、640 像素、FP16、融合卷积、前向推理加 NMS、80 次预热和 300 次计时迭代。候选模型与 YOLOv8n 在同一进程和同一 GPU 中加载，并按 AB/BA 顺序交替测试。本文同时报告绝对延迟、FPS 和对应的同场基线结果。
